\section{Output Analysis}

\subsection{Design}

Three TO-BE models were made: an optimised patient schedule with even appointment times, a pre-processing model where residents review cases the day before, and a combined scenario applying both these modifications.

The Miller Pain Treatment Center is a terminating system, with each session starting empty and ending when all patients have checket out of the system. This means there is no need for a warm-up.
As each session lasts approximately 4 hours, we set the simulation time to 6 hours to ensure that every patient could finish their visit. All models where run with 30 replications, as the confidence intervals of all replications were below 8 minutes. The random seeds used in the probability distributions were consistent for all models, enabling comparison and reducing variance accross simulations.

\subsection{Performance measures}

The primary outout variable, as described by Dr. Weem, is mean cycle time per patient \parencite{MillerPainTreatmentCenter}. This is defined as time in minutes from arrvial to checkout. Patients processed is a secondary concern, meaning the number of patients processed through the system for each replication.

\subsection{Confidence intervals}

\autoref{tab:ci} presents the mean cycle time and 95\% significance intervals for all models across 30 replications. The results shows a clear improvement of the system for all modifications, with the combined solution giving the best performance. The reduction in confidence interval across all models show not only an improvement in performance, but also a reduction in variability, suggesting a more consistent and reliable process flow.

\begin{table}[h]
\centering
\begin{threeparttable}
\begin{tabular}{lrrrr}
\toprule
\textbf{Model} & \textbf{Mean} & \textbf{Std} & \textbf{CI lower} & \textbf{CI upper} \\
\midrule
AS-IS                   & $73.95$ & $20.81$ & $66.18$ & $81.72$ \\
TO-BE 1: Scheduling     & $65.60$ & $15.34$ & $59.87$ & $71.33$ \\
TO-BE 2: Pre-processing & $60.99$ & $14.80$ & $55.46$ & $66.51$ \\
TO-BE 3: Combined       & $53.95$ & $10.31$ & $50.10$ & $57.80$ \\
\bottomrule
\end{tabular}
\begin{tablenotes}
\small
\item 
\end{tablenotes}
\end{threeparttable}
\caption{Mean cycle time and 95\% confidence intervals across models}
\label{tab:ci}
\end{table}

\subsection{Statistical comparison}

To compare the output variables from the different scenarios we perform a paired t-test. As all models use the same random seeds, each scenario is exposed to the same set of stochastic inputs, meaning the difference in performance comes from modifications done to the TO-BE models rather than randomness. The difference is defined as $D_i = \bar{x}_i^{TO-BE} - \bar{x}_i^{AS-IS}$, with the hypothesis tested with a significance below 0.05.

As each replication produces a mean cycle time, 30 iterations provides 30 independent observations of this variable. Since the values vary due to randomness, we make a 95\% confidence interval for the sample mean. To calculate this interval we use t-distribution with $n - 1 = 29$ degrees of freedom, as the standard deviation in the population is unkown and therefore estimated from the data. The interval is calculated as $\bar{x} \pm t_{0.025,29} \cdot \frac{s}{\sqrt{30}}$, where $x$ is the mean cycle time, $s$ is the standard deviation and $t_{0.025,29}$ is the critical value. 

\begin{table}[H]
\centering
\begin{threeparttable}
\begin{tabular}{lrrrrr}
\toprule
\textbf{Comparison} & \textbf{$\bar{D}$} & \textbf{CI lower} & \textbf{CI upper} & \textbf{t-statistic} & \textbf{p-value} \\
\midrule
Scheduling vs.\ AS-IS      & $-6.71$  & $-8.12$  & $-5.31$ & $-9.80$  & $<0.001$ \\
Pre-processing vs.\ AS-IS  & $-5.47$  & $-6.76$  & $-4.19$ & $-8.71$  & $<0.001$ \\
Combined vs.\ AS-IS        & $-12.11$ & $-14.41$ & $-9.81$ & $-10.76$ & $<0.001$ \\
\bottomrule
\end{tabular}
\end{threeparttable}
\caption{Paired t-test results}
\label{tab:ttest}
\end{table}